\section{Ping-pong and the Tits alternative}\label{sec:tits}

\begin{theorem}[Tits~\cite{tits-1972}]\label{thm:tits}
Let $K$ be a field and let $\Gamma \le \GL_n(K)$ be a finitely generated
subgroup. Then either $\Gamma$ contains a free subgroup of rank~$2$, or
$\Gamma$ has a solvable subgroup of finite index.
\end{theorem}

The two cases exclude each other: a subgroup of a virtually solvable group is
virtually solvable, and a free group of rank~$2$ is not. The free subgroup in
Tits' proof comes from ping-pong between high powers of two well-chosen
elements. Here is the simplest instance, with a complete proof.

\begin{proposition}\label{prop:hyperbolic}
Let $g, h \in \SL_2(\R)$ with $\abs{\tr(g)} > 2$ and $\abs{\tr(h)} > 2$, and
suppose that $g$ and $h$ have no common fixed point in $\RP^1$. Then there is
$N \ge 1$ such that $g^N$ and $h^N$ freely generate a free group of rank~$2$.
\end{proposition}

\begin{proof}
Since $\abs{\tr(g)} > 2$, the matrix $g$ has real eigenvalues $\lambda$ and
$\lambda^{-1}$ with $\abs{\lambda} > 1$. Their eigenlines are the two fixed
points of $g$ in $\RP^1$: the attracting point $g^+$ (for $\lambda$) and the
repelling point $g^-$ (for $\lambda^{-1}$).

\emph{Claim.} For all open sets $U^+ \ni g^+$ and $U^- \ni g^-$ there is $N_0$
such that, for every $n \ge N_0$,
\[
  g^n(\RP^1 \setminus U^-) \subseteq U^+ \quad\text{and}\quad
  g^{-n}(\RP^1 \setminus U^+) \subseteq U^-.
\]
In a basis of eigenvectors $g = \mathrm{diag}(\lambda, \lambda^{-1})$. Give
the line through $(x,y)$ the coordinate $z = x/y \in \R \cup \{\infty\}$, so
that $g^+ = \infty$, $g^- = 0$, and $g^n$ acts by $z \mapsto \lambda^{2n} z$.
The compact set $\RP^1 \setminus U^-$ avoids $0$, so $\abs{z} \ge \delta$ on
it for some $\delta > 0$, and $U^+$ contains $\{\abs{z} > R\}$ for some $R$.
Once $\lambda^{2n}\delta > R$, the first inclusion holds. The second is the
same argument for $g^{-1}$, whose attracting point is $g^-$.

The four points $g^\pm, h^\pm$ are distinct, so they have pairwise disjoint
open neighborhoods $U_g^\pm$ and $U_h^\pm$. Choose $N$ so that the claim
holds for $g$ and for $h$ with these neighborhoods, and set $a = g^N$,
$b = h^N$ and
\[
  D_a = U_g^+, \quad D_{a^{-1}} = U_g^-, \quad D_b = U_h^+, \quad D_{b^{-1}} = U_h^-.
\]
If $t \ne a^{-1}$, then $D_t \subseteq \RP^1 \setminus U_g^-$, so
$a(D_t) \subseteq D_a$. The other three letters work the same way, and
Theorem~\ref{thm:pingpong} applies.
\end{proof}

\begin{proof}[Sketch of the general case]
This is an outline, not a proof; the details are the heart
of~\cite{tits-1972}. Let $k$ be a local field, such as $\R$, $\C$ or the
$p$-adic numbers $\Q_p$. Call $g \in \GL_n(k)$ \emph{proximal} if it has a
unique eigenvalue of largest absolute value, and that eigenvalue is simple.
On the projective space $\mathbb{P}(k^n)$, high powers of $g$ push everything
that stays away from a repelling hyperplane (the sum of the other generalized
eigenspaces) close to an attracting point (the eigenline of that
eigenvalue). Hyperbolic elements of $\SL_2(\R)$ are the case $n = 2$.
\begin{enumerate}
\item \emph{Ping-pong.} If $g^{\pm1}$ and $h^{\pm1}$ are all proximal, and each
attracting point lies off the repelling hyperplanes of the other element and
its inverse, then the proof of Proposition~\ref{prop:hyperbolic}, with small
neighborhoods of the four attracting points as the sets $D_s$, shows that
$g^N$ and $h^N$ freely generate a free group of rank~$2$ for large $N$.
\item \emph{Finding such elements.} If $\Gamma$ is not virtually solvable, one
reduces to the case where its Zariski closure is semisimple. One then embeds
the field generated by the matrix entries of $\Gamma$ into a local field $k$,
chosen so that some $\gamma \in \Gamma$ has eigenvalues of different absolute
values. A suitable representation, such as an exterior power, makes $\gamma$
proximal, and Zariski density supplies a conjugate of $\gamma$ in general
position with respect to $\gamma$. \qedhere
\end{enumerate}
\end{proof}

The choice of local field matters: a compact group such as $\SO_3(\R)$ acts
on the real projective plane by isometries, so the contraction has to be found
over another local field, for example a $p$-adic one. Breuillard and Gelander~\cite{breuillard-gelander-2007} proved a
topological version of the alternative, in which ping-pong on projective
spaces over local fields produces free subgroups that are dense, for example
in compact Lie groups.
